\hmsubsection{Improvements}{}{}\label{sec:improvements}

This section identifies concrete improvements that would
address the limitations described in Section~\ref{sec:limitations} and extend
the system's capability beyond its current state. The
suggestions are ordered by estimated impact relative to
implementation effort.

\hmsubsubsection{Kinematic and Hardware Improvements}{UMA}{}

\textbf{Reduce $L_3$ assembly mass.} A remaining hardware improvement
for the present arm would be to further reduce the weight of the
wrist-to-claw assembly. The current $L_3$ link
(22.0~cm) carries the OAK-D~S2 camera, the camera mount, and
the gripper, which together create a significant gravitational
torque on the elbow motor in unsupported poses. During integration,
reducing the mass of these components improved the thermal situation
somewhat, and the larger practical improvement for idle scanning came
from revising the scan pose so that the $L_2$ structure rests on the
$L_1$ piece. Further mass reduction would therefore mainly improve
dynamic margins, reduce sag, and provide additional robustness during
movement and outside the mechanically supported scan posture.

\textbf{Lighter camera.} The OAK-D~S2 provides depth sensing
and onboard neural inference, but the present application uses
only the RGB stream for HSV-based ball detection. The depth
and inference capabilities are not utilised, making the
OAK-D~S2 heavier and more expensive than necessary. A lighter
alternative such as the OAK-D~Lite, which provides a
comparable RGB sensor in a smaller and lighter form factor,
would reduce the mass of the $L_3$ assembly and thereby lower
the gravitational torque on the elbow. Since the deployed scan pose now
mechanically supports the $L_2$ structure on the $L_1$ piece, a lighter
camera would primarily add thermal and positioning margin for movements
and any unsupported postures rather than serving as the only mitigation
for idle scanning.

\textbf{Wrist motor upgrade.} The wrist joint (M4) currently
uses an XL430-W250-T, which has a stall torque of 1.5~N$\cdot$m at 12~V.
Replacing it with an XM430-W350-T (4.1~N$\cdot$m stall torque)
would increase the wrist's holding capacity and reduce the risk
of droop at extended reach. The XM430 series also provides the
present current register used by the adaptive grip algorithm,
which is not available on the XL430. Unifying all five motors
to the same family would simplify the spare-part inventory and
the firmware's error-handling logic. Since movement-related heating
remained observable even after idle scan-pose heating was mitigated,
higher actuator margins would also reduce thermal risk during repeated
sorting cycles.

\textbf{Elbow motor torque upgrade.} The elbow motor (M3)
currently uses a Dynamixel XM430-W210-T. Replacing it with the
XM430-W350-T, which has the same physical dimensions and mass
but a higher stall torque (4.1~N$\cdot$m versus 3.0~N$\cdot$m),
would increase the torque margin during unsupported movements and
reduce the motor's operating point as a fraction of its rated capacity.
A motor that operates further below its stall torque draws less
current and generates less heat, which would add robustness to both the
idle scan-pose mitigation described in Section~\ref{sec:hardware-integration} and unsupported
movement phases. This upgrade was
considered during the project but was not pursued due to
uncertainty about whether additional funding would be approved
for motor procurement.

\textbf{Physical sag model.} The current sag compensation uses
a phenomenological polynomial fit that captures the observed
droop but does not model the underlying mechanics. A
physics-based model that accounts for the arm's mass
distribution, joint compliance, and motor backlash would
generalise better across the workspace and reduce the need
for dense calibration points. Such a model would also enable
sag prediction for configurations that were not sampled
during calibration.

\textbf{Force-sensing gripper.} The adaptive grip algorithm
uses the Dynamixel motor's built-in current and load registers
as a proxy for grip force. A dedicated force sensor
(for example, a strain gauge or force-sensitive resistor) in
the gripper jaws would provide higher-resolution feedback and
would allow the system to measure the actual grip force in
Newtons rather than inferring it from motor current. This
would improve reliability for objects of different sizes and
materials.

\textbf{Six-degree-of-freedom arm.} Upgrading to a
six-degree-of-freedom kinematic chain would remove the
constraint that the gripper must point downwards during
pickup. This would expand the reachable workspace and allow
the arm to approach objects from arbitrary angles, but would
require replacing the geometric solver with a numerical or
analytical method capable of handling the additional degrees
of freedom.
